\section{Introduction}

Single-photon LiDAR enables time-resolved optical ranging in conditions where only a small number of returned photons are available. Such conditions arise at long range, for low-reflectivity objects, under atmospheric attenuation, and when the transmitted optical energy must be constrained.

Most conventional acquisition strategies distribute measurements uniformly across the sensing field. Uniform allocation is simple and reproducible, but it may devote substantial photon resources to regions that have already been rejected or provide little information. Adaptive sensing instead updates the acquisition strategy after each observation and concentrates future measurements on regions expected to improve the target-presence decision.

This work investigates physics-guided adaptive photon allocation for defense-oriented single-photon LiDAR reconnaissance. The proposed strategy combines Bayesian target uncertainty with an optical signal-to-background prediction derived from a physically grounded LiDAR model. The allocation rule also includes an exploration term that prevents repeated concentration on a single initially favored region.

The principal contributions are:

\begin{enumerate}
    \item A sequential photon-allocation framework incorporating atmospheric attenuation, distributed backscatter, beam divergence, target-footprint coupling, and SPAD detector nonidealities.
    \item A physics-guided information-allocation strategy combining posterior uncertainty, predicted optical contrast, and measurement-count exploration.
    \item A reproducible comparison with uniform, random, greedy-posterior, and entropy-based allocation under a common photon budget.
    \item A quantitative demonstration of improved detection reliability and reduced average photon consumption relative to uniform sensing.
\end{enumerate}
